\chapter*{Anexos}
En este anexo se detallan las instrucciones necesarias para acceder al código, así como los requisitos técnicos para compilar y ejecutar el simulador gráfico de incendios forestales de forma local.

\section*{Acceso al Código Fuente}
El proyecto está versionado utilizando el sistema Git y alojado en la plataforma GitHub. El repositorio incluye tanto el código fuente de Godot (GDScript) como los \textit{Shaders} de la GPU para el correcto funcionamiento del Autómata Celular y el \textit{Raymarching}.

\begin{itemize}
    \item \textbf{Enlace al repositorio:} \url{https://github.com/GerardoBonet25/TFG}
    \item \textbf{Rama principal:} \texttt{main} (contiene la versión estable documentada en esta memoria).
\end{itemize}
En este directorio se encuentran los archivos tanto del simulador principal (SimuladorIncendios) como los archivos de la memoria y un pequeño simulador inicial 2D que realicé para comprender las bases (InicioPropagacion2D).\\

\section*{Requisitos del Sistema}
Debido al uso de tecnologías de renderizado avanzado y paralelización masiva en GPU, el equipo informático destinado a ejecutar la simulación debe cumplir con los siguientes requisitos mínimos:\\

\begin{itemize}
    \item \textbf{Sistema Operativo:} Windows 10/11.
    \item \textbf{Motor Gráfico:} Godot Engine versión 4.2 o superior.
    \item \textbf{API Gráfica:} Tarjeta gráfica (GPU) dedicada o integrada con soporte completo para \textbf{Vulkan 1.2}. Este requisito es fundamental ya que la API utilizada para despachar los \textit{Compute Shaders} del Autómata Celular no es compatible con versiones antiguas de OpenGL.

\end{itemize}

